\chapter{Findings}
\label{chap:findings}

\section{Numerical summary}
Method B was higher by \SI{0.18}{\milli\metre} on average in the simulated
sample (\cref{tab:findings}).

\begin{table}[htbp]
  \centering
  \caption{Decision-relevant summary of simulated results.}
  \label{tab:findings}
  \begin{tabularx}{\linewidth}{@{}
    >{\raggedright\arraybackslash}X
    S[table-format=2.2]
    S[table-format=1.2]
    >{\raggedright\arraybackslash}X@{}}
    \toprule
    Quantity & {Estimate (\si{\milli\metre})}
      & {SD (\si{\milli\metre})} & Interpretation \\
    \midrule
    Method A & 12.44 & 1.31 & Reference values \\
    Method B & 12.62 & 1.35 & Candidate values \\
    Difference & 0.18 & 0.24 & Positive offset \\
    \bottomrule
  \end{tabularx}
\end{table}

\section{Visual assessment}
\begin{figure}[htbp]
  \centering
  \includegraphics[width=.8\textwidth]{difference_plot.pdf}
  \caption{Simulated paired differences across the observed range. The labelled
  reference lines are descriptive, not acceptance thresholds.}
  \label{fig:differences}
\end{figure}

The display does not by itself show whether the variability is tolerable. That
judgement belongs to a documented decision rule.
